% ACC 2027 compact paper structure.
%
% IMPORTANT:
% 1. Obtain the current ieeeconf.cls and sample file from the ACC 2027
%    PaperPlaza Support/Author Kit and place ieeeconf.cls on the TeX path.
% 2. ACC 2027 requires the standard two-column format. Initial submissions are
%    limited to 8 pages; accepted papers are 6 pages at no charge and may use
%    up to 2 paid extra pages.
% 3. Replace all TODO fields and reconcile every result with
%    docs/FULL_PROJECT_AUDIT_2026-09-02.md before submission.
% 4. ACC 2027 requires transparent disclosure and citation of LLM use.

\documentclass[letterpaper,10pt,conference]{ieeeconf}

\IEEEoverridecommandlockouts
\overrideIEEEmargins

\usepackage{amsmath,amssymb,bm}
\usepackage{booktabs,multirow,array}
\usepackage{graphicx}
\usepackage{xcolor}
\usepackage{cite}
\usepackage[hidelinks]{hyperref}

\graphicspath{{../results/reports/}{../results/plots/}{../presentation/figures/}}

\newcommand{\R}{\mathbb{R}}
\newcommand{\T}{^{\mathsf T}}
\newcommand{\sat}{\operatorname{sat}}
\newcommand{\clip}{\operatorname{clip}}
\newcommand{\sgn}{\operatorname{sgn}}

\title{\LARGE \bf
Online PES-Adapted Neuromorphic Membrane--Gate Residual Control\\
for Actuator-Constrained UAV Quaternion Attitude Tracking
}

\author{Gandhi A. Contreras Torres%
\thanks{Author affiliation, address, funding, and correspondence information
must be completed before submission.}%
}

\begin{document}

\maketitle
\thispagestyle{empty}
\pagestyle{empty}

\begin{abstract}
% TODO: Keep the abstract to one compact paragraph.
% Recommended order:
% (1) actuator-constrained quaternion-attitude problem;
% (2) classical PD plus bounded additive CMG residual;
% (3) paired membrane--gate dynamics and online PES decoder;
% (4) primary X2 and explicitly labeled CF2 1.5x design results;
% (5) repeated physical Hover Roll result and mission-dependent hardware matrix;
% (6) principal limitation: no completed full-scale physical yaw360.
This paper develops a baseline-preserving neuromorphic residual controller for
actuator-constrained quaternion attitude tracking. A classical proportional--
derivative controller supplies the stabilizing body-torque command, while
paired continuous coupled membrane--gate units and an online PES-adapted
decoder generate a bounded additive correction. The method is implemented
both directly and as a custom non-spiking Nengo neuron type, with explicit
rotor allocation and complete neural-state telemetry. Matched MuJoCo and
Quanser 3DOF Hover experiments characterize tracking improvement, residual
authority, actuator feasibility, and the limits of simulation-to-hardware
transfer.
\end{abstract}

% ACC does not require a separate keyword block in the standard sample.

\section{Introduction}
\label{sec:introduction}

% Compact narrative:
% - Why actuator constraints make direct neural replacement problematic.
% - Why the baseline must be preserved.
% - Why a continuous neuromorphic dynamical state is useful.
% - Gap: online local adaptation plus explicit rotor allocation and hardware.
% - End with a numbered contribution list.

\subsection{Motivation and Research Gap}

% Introduce learning-enabled attitude control, transparency, actuator
% feasibility, and simulation-to-hardware transfer.

\subsection{Contributions and Claim Boundaries}

% Recommended contributions:
% 1. Continuous CMG dynamics with paired directional encoding.
% 2. Online Nengo/PES decoder in a bounded additive residual architecture.
% 3. Explicit rotor/voltage allocation and stock-CF2 authority audit.
% 4. Direct-versus-Nengo equivalence plus MuJoCo mission matrix.
% 5. Physical 3DOF Hover validation with repeated Roll evidence.
% State explicitly: no universal superiority over all neuron models and no
% completed full-scale physical yaw360.

\section{Problem Formulation}
\label{sec:problem}

\subsection{Rigid-Body Attitude Dynamics}

Let $q\in\mathbb{S}^{3}$ be the body attitude quaternion,
$\bm\omega\in\R^3$ the body angular velocity, $\bm J$ the inertia tensor,
and $\bm\tau$ the applied body torque. The rotational dynamics are
\begin{align}
    \dot q &= \frac{1}{2}q\otimes[0,\bm\omega\T]\T, \\
    \bm J\dot{\bm\omega}
    &= \bm\tau-\bm\omega\times(\bm J\bm\omega).
\end{align}

\subsection{Quaternion Error and Tracking Objective}

For desired attitude $q_d$, define
\begin{align}
    q_e &= q_d\otimes q^{-1}, \\
    \bm e_q &= \sgn(q_{e,0})\operatorname{vec}(q_e), \\
    \theta_q &= 2\cos^{-1}\!\left(\clip(|q_{e,0}|,0,1)\right).
\end{align}

% Explain shortest-path sign selection and distinguish quaternion-angle error
% from the wrapped Euler-component norm used in physical Hover summaries.

\subsection{Actuator-Constrained Control Objective}

% Define bounded control channels u0, limits ell, nonnegative rotor commands,
% and the tracking/authority/saturation objectives.

\section{Baseline Quaternion Attitude Control}
\label{sec:baseline}

\subsection{Quaternion PD Law}

With natural frequencies $\bm\omega_n$ and damping ratio $\zeta$,
\begin{align}
    \bm K_q &= 2\operatorname{diag}(\bm J)
        \odot\bm\omega_n^2, \\
    \bm K_{\omega} &= 2\zeta\operatorname{diag}(\bm J)
        \odot\bm\omega_n, \\
    \bm\tau_0 &= \sat_{\bm\tau_{\max}}
        \left(-\bm K_q\odot\bm e_q
        -\bm K_{\omega}\odot\bm\omega\right).
\end{align}

\subsection{Collective Thrust and Rotor Allocation}

For four rotor thrusts $\bm f\in\R_{geq0}^{4}$,
\begin{equation}
    \begin{bmatrix}T&\tau_x&\tau_y&\tau_z\end{bmatrix}\T
    =\bm B\bm f,
    \qquad
    \bm f=\bm B^{-1}
    \begin{bmatrix}T&\bm\tau\T\end{bmatrix}\T.
\end{equation}

% Define B from rotor positions and reaction-torque coefficients. State that
% collective thrust remains classical and CMG augments only body torque.

\section{Online PES-Adapted Membrane--Gate Residual Control in Nengo}
\label{sec:cmg_nengo}

\subsection{Paired Directional Current Encoding}

For $m$ bounded baseline channels $\bm u_0$ with symmetric limits
$\bm\ell>0$, define
\begin{align}
    \bar{\bm u} &= \bm u_0\oslash\bm\ell, \\
    \bm I^+ &= I_b\bm 1+I_g\max(\bar{\bm u},\bm 0), \\
    \bm I^- &= I_b\bm 1+I_g\max(-\bar{\bm u},\bm 0).
\end{align}

% UAV attitude: m=3 and 2m=6 units. Explain positive/negative channel order.

\subsection{Coupled Membrane--Gate Dynamics}

For unit $j$,
\begin{align}
    \dot r_j &= \frac{\sigma\!\left(
        \alpha_h[v_j-(v_{\mathrm{th}}-\gamma r_j)]\right)-r_j}
        {\tau_r}, \\
    \dot v_j &= -\frac{v_j-v_{\mathrm{rest}}}{\tau}
        +\frac{I_j}{\tau}
        -c_j(v_j-v_{\mathrm{reset}}), \\
    c_j &= \left(g_0+\mathbf{1}_{\mathrm{adapt}}g_1|I_j|\right)r_j^2.
\end{align}

% Make the conceptual distinction explicit:
% - v and r are mutually coupled internal states;
% - the decoder reads gate-derived activity, not v independently;
% - this implementation is continuous and non-spiking.

\subsection{Continuous Activity and Fixed Decoder}

\begin{align}
    \bm a_k &= \clip_{[0,1]}(G\bm r_k), \\
    \bm W_0 &= \begin{bmatrix}\bm I_m&-\bm I_m\end{bmatrix}, \\
    \bm s_k &= \clip_{[-1,1]}(\bm W_k\bm a_k).
\end{align}

\subsection{Online PES Decoder Adaptation}

\begin{align}
    \bm y_k^* &= \clip_{[-1,1]}
        (-k_e\bm e_{q,k}-k_\omega^*\bm\omega_k), \\
    \bm\varepsilon_k &= \bm y_k^*-\bm s_k, \\
    \bm W_{k+1} &= \Pi_{[-W_{\max},W_{\max}]}
    \left(\bm W_k+\eta\bm\varepsilon_k\bm a_k\T
    -\lambda_W\bm W_k\right).
\end{align}

% Explain causality: s_k is decoded with W_k; W_{k+1} affects the next sample.

\subsection{Bounded Additive Residual}

\begin{align}
    \bm u_k^{\mathrm{neural}} &= \bm s_k\odot\bm\ell, \\
    \bm u_k &= \sat_{\bm\ell}
        \left(\bm u_{0,k}+\beta\bm u_k^{\mathrm{neural}}\right), \\
    \bm\delta_k &= \bm u_k-\bm u_{0,k}.
\end{align}

% Contrast with the historical noncanonical blend
% (1-beta)u0 + beta*u_neural.

\subsection{Adaptive Conductance Versus Adaptive Decoder}

% State the audited result: fixed/adaptive conductance CMG variants are
% numerically indistinguishable in the amplified-gate regime. The large gain
% comes from online W adaptation, not from g1|I|.

\subsection{Nengo Realization and Numerical Equivalence}

% Describe CoupledMembraneGate as a custom nengo.neurons.NeuronType, direct
% current injection, Nengo built-in PES, learning-rate conversion, weight
% projection/leak, and the approximately 1.05e-7 isolated relative command
% error versus the direct NumPy implementation.

\subsection{Timing and Neural Telemetry}

% MuJoCo: 2 kHz plant/neural step. Hover: 200 Hz Q8, ten 0.5-ms Nengo steps,
% 2 kHz neural archive, optional 50 Hz GUI. List v, r, activity, current,
% conductance, threshold, decoder output, PES error, W, and Delta W.

\section{Platform Realization and Experimental Methodology}
\label{sec:platform_method}

\subsection{Skydio X2 Rotor-Actuated MuJoCo Model}

% Identify the Menagerie model, four named rotor actuators, mixer, bounds,
% vertical loop, and saturation measurement.

\subsection{Crazyflie Rotor-Authority Design and Stock Boundary}

% Separate two cases:
% (1) stock Menagerie CF2 abstract moment interface: feasibility boundary;
% (2) explicit four-rotor design with 1.5x nominal rotor authority: design study.
% Never call the second case stock CF2 hardware.

\subsection{Quanser 3DOF Hover Hardware Interface}

% Q8 encoder and actuator ordering, startup tare, quaternion reconstruction,
% filtered body rates, torque-to-F/B/L/R allocation, 0--24 V command,
% amplifier gain 3, residual cap/slew, and unconditional zero-output cleanup.

\subsection{Reference Missions}

% Define Hover, Roll, Full3, Hard45, Hardstep, and Yaw360 mathematically or in
% a compact table. Record scale and duration independently.

\subsection{Matched Controller Protocol}

% Same reference, duration, scale, allocator, initial-condition convention,
% safety limits, and logging for baseline/CMG pairs. State n for every table.

\subsection{Parameters and Evaluation Metrics}

% Include one compact parameter table. Distinguish:
% - MuJoCo quaternion-angle mean/max [rad];
% - Hover wrapped Euler-vector mean/max [deg];
% - normalized residual RMS/peak;
% - decoder drift/update;
% - actuator saturation and mission completion.

\subsection{Hardware Timing and Safety Protocol}

% Explicit arm and confirmation gates; qualified controllers only; encoder
% step, pitch/roll, deadline, voltage, residual, and cleanup guards.

\section{MuJoCo UAV Results}
\label{sec:mujoco_results}

\subsection{Five-Mission Quaternion-Attitude Results}

% Primary table: X2 and CF2 1.5x design, Roll/Full3/Hard45/Hardstep/Yaw360.
% Report PD, fixed CMG, CMG--PES, improvement, residual authority, saturation.

\subsection{Decoder Adaptation and Conductance Ablation}

% Show fixed decoder versus PES-adapted decoder and fixed versus adaptive
% conductance. Attribute the large improvement only to decoder adaptation.

\subsection{Actuator Utilization and Feasibility}

% Include X2 zero-saturation evidence, stock-CF2 failure boundary, and
% explicit labeling of the 1.5x design case.

\subsection{Six-DOF Extension}

% Use one compact table for X2 waypoint, circle, and full6-stress position and
% attitude metrics. Keep ideal-wrench and real-rotor results separate.

\section{Neuromorphic Model Benchmark}
\label{sec:neuron_benchmark}

\subsection{Comparator Models and Matched Decoder Conditions}

% CMG, LIF, ALIF, Izhikevich, and LIFRate; fixed and PES decoders. Explain
% feature calibration and why sparse activity is an operating-point result.

\subsection{Tracking Performance and Residual Authority}

% Report that Izhikevich--PES can win individual cells, while CMG--PES uses
% lower normalized residual RMS and decoder movement. Avoid universal claims.

\subsection{Interpretation and Benchmark Limitations}

% Discuss seed count, matched current encoding, neuron-specific state/readout,
% absence of armed non-CMG hardware trials, and need for frozen-weight tests.

\section{Physical 3DOF Hover Results}
\label{sec:hover_results}

\subsection{Direct-Nengo Execution and Data Integrity}

% 200 Hz Q8 worker, 2 kHz Nengo, full telemetry, 12 canonical completed runs,
% 144/144 deep checks. Mention why the GUI was removed from the hard loop.

\subsection{Repeated Roll Experiment}

% Primary physical result: n=3 matched pairs, PD 5.4405+/-0.0648 deg,
% CMG 5.2960+/-0.0428 deg, paired mean improvement 2.6508+/-1.0765%.

\subsection{Complete Matched Mission Matrix}

% Report all valid rows, including negative results:
% Hover +60.38% near zero; Roll +5.57%; Full3 -3.21%; Hard45 +0.55%;
% Hardstep +0.95%; reduced-range Yaw -8.38%.
% Include maximum-error changes, which are negative in several cells.

\subsection{Yaw Feasibility Boundary}

% The completed scale-0.3 profile ends near 108 deg and is not yaw360.
% Two full-scale PD attempts aborted at approximately +/-20.127 deg roll.
% Present mechanical/orientation/cable wind-up only as a hypothesis.

\section{Discussion}
\label{sec:discussion}

\subsection{Why Online PES Produces the Main Improvement}

% Link W adaptation, tracking target, phase-lag correction, and mission demand.

\subsection{Roles of Membrane State and Gate State}

% The decoder reads gate-derived activity, but v and r are coupled and both
% participate in the internal dynamics. Do not state that r acts alone.

\subsection{Performance--Authority and Actuator Tradeoffs}

% Compare simulation beta/authority with the deliberately conservative Hover
% residual cap and slew. Discuss mean versus maximum-error tradeoffs.

\subsection{Simulation-to-Hardware Transfer and Limitations}

% Single-run mission cells, hardware variation, reduced reference scales,
% model uncertainty, absent full yaw360, online-on-evaluation adaptation,
% lack of frozen-weight generalization, and non-CMG hardware limitation.

% If page space permits, mention the QUBE QLabs result in one paragraph:
% stable 500-Hz transfer and practical equivalence to LQR, not improvement.
% Keep QCar and QDrone as future work, not completed evidence.

\section{Conclusion}
\label{sec:conclusion}

% Summarize only audited claims. End with repeated physical trials,
% disturbance/held-out/frozen-weight evaluation, and neuromorphic perception
% integration as future work.

\section*{Acknowledgment}

% TODO: Add advisor, laboratory, funding, equipment, and contributor credits.
% TODO: Add the ACC 2027-required transparent LLM-use disclosure and cite the
% tools in the reference list. Describe exactly which research, coding,
% analysis, visualization, and writing tasks used each tool. The human authors
% retain responsibility for validation and all claims.

\begin{thebibliography}{99}

% TODO: Replace placeholders with complete references used in the manuscript.
% Use \bibitem entries or the ACC-approved BibTeX workflow from the author kit.

\bibitem{bekolay2013pes}
A.~Bekolay et al., ``Nengo: A Python tool for building large-scale functional
brain models,'' \emph{Frontiers in Neuroinformatics}, vol.~7, 2014.

% TODO: Insert verified quaternion attitude-control, PES, multirotor,
% actuator-allocation, and neuromorphic-control references.

\end{thebibliography}

\end{document}
